\documentclass[11pt]{article}
\usepackage[utf8]{inputenc}
\usepackage{amsmath,amssymb}
\usepackage{hyperref}
\title{Robust Continual Learning Catastrophic Forgetting under Distribution Shift}
\author{Anonymous}
\date{}
\begin{document}
\maketitle

\begin{abstract}
This paper studies Continual Learning Catastrophic Forgetting. Grounded in a real, citable literature \cite{ref1,ref2,ref3,ref4,ref5,ref6,ref7,ref8},
we propose a focused method and report \textbf{illustrative} experiments.
All references below are real and verifiable (OpenAlex/arXiv); the reported
numbers are simulated to illustrate intended behaviour.
\end{abstract}

\section{Introduction}
Continual Learning Catastrophic Forgetting is an active research area. This work targets a concrete gap identified from
the recent literature and contributes a method together with an evaluation protocol.

\section{Related Work (real literature)}
The following works, retrieved from public bibliographic APIs, frame our study:
\begin{itemize}
\item Parisi et al. (2019) — "Continual lifelong learning with neural networks: A review", Neural Networks (cited 3060).
\item Delange et al. (2021) — "A continual learning survey: Defying forgetting in classification tasks", IEEE Transactions on Pattern Analysis and Machine Intelligence (cited 1611).
\item Wang et al. (2024) — "A Comprehensive Survey of Continual Learning: Theory, Method and Application", IEEE Transactions on Pattern Analysis and Machine Intelligence (cited 847).
\item Wang et al. (2022) — "Learning to Prompt for Continual Learning", 2022 IEEE/CVF Conference on Computer Vision and Pattern Recognition (CVPR) (cited 675).
\item Ven et al. (2019) — "Three scenarios for continual learning", Lirias (cited 557).
\item Zhao et al. (2020) — "Maintaining Discrimination and Fairness in Class Incremental Learning" (cited 475).
\end{itemize}

\section{Method}
We model distribution shift explicitly and optimise a worst-case objective over an uncertainty set of plausible test distributions. We instantiate this idea for Continual Learning Catastrophic Forgetting and analyse its cost/quality trade-off.

\section{Experiments (Illustrative)}
\begin{center}
\begin{tabular}{lcc}
System & Primary Metric & Efficiency \\ \hline
Strong baseline & 0.812 & 1.00$\times$ \\
Ablation & 0.828 & 0.92$\times$ \\
\textbf{Ours} & \textbf{0.851} & \textbf{0.71$\times$}
\end{tabular}
\end{center}
\emph{Metrics simulated to illustrate the intended effect; not measured.}

\section{Conclusion}
We connected a real literature to a concrete method for Continual Learning Catastrophic Forgetting. Future work will
validate the illustrative results with real experiments.

\begin{thebibliography}{9}
\bibitem{ref1} German I. Parisi, Ronald Kemker, Jose L. Part et al.. Continual lifelong learning with neural networks: A review. \emph{Neural Networks}, 2019. DOI:10.1016/j.neunet.2019.01.012
\bibitem{ref2} Matthias Delange, Rahaf Aljundi, Marc Masana et al.. A continual learning survey: Defying forgetting in classification tasks. \emph{IEEE Transactions on Pattern Analysis and Machine Intelligence}, 2021. DOI:10.1109/tpami.2021.3057446
\bibitem{ref3} Liyuan Wang, Xingxing Zhang, Hang Su et al.. A Comprehensive Survey of Continual Learning: Theory, Method and Application. \emph{IEEE Transactions on Pattern Analysis and Machine Intelligence}, 2024. DOI:10.1109/tpami.2024.3367329
\bibitem{ref4} Zifeng Wang, Zizhao Zhang, Chen‐Yu Lee et al.. Learning to Prompt for Continual Learning. \emph{2022 IEEE/CVF Conference on Computer Vision and Pattern Recognition (CVPR)}, 2022. DOI:10.1109/cvpr52688.2022.00024
\bibitem{ref5} Gido M. van de Ven, Andreas S. Tolias. Three scenarios for continual learning. \emph{Lirias}, 2019. DOI:10.48550/arxiv.1904.07734
\bibitem{ref6} Bowen Zhao, Xi Xiao, Guojun Gan et al.. Maintaining Discrimination and Fairness in Class Incremental Learning. \emph{}, 2020. DOI:10.1109/cvpr42600.2020.01322
\bibitem{ref7} Gido M. van de Ven, Hava T. Siegelmann, Andreas S. Tolias. Brain-inspired replay for continual learning with artificial neural networks. \emph{Nature Communications}, 2020. DOI:10.1038/s41467-020-17866-2
\bibitem{ref8} Zifeng Wang, Zizhao Zhang, Sayna Ebrahimi et al.. DualPrompt: Complementary Prompting for Rehearsal-Free Continual Learning. \emph{Lecture notes in computer science}, 2022. DOI:10.1007/978-3-031-19809-0\_36
\end{thebibliography}
\end{document}
